Emmy stond in de kou.
Bloot was niet wat ze wou,

maar ze was geen vechter.
“Waarom wordt het slechter

terwijl het al slecht is
en de dood ons gewis

het slechtste zal geven?”.
Emmy bleef maar beven.

Voltaire’s Candide
beschreef haar timide

pessimistisch idee.
Leibniz’ theodicee

over Gods erbarmen,
kon ze niet omarmen.

Schrik had zich gevestigd.
Ze werd overweldigd

door alle sensaties.
Angst baarde emoties

als paniek en woede.
Ze was op haar hoede,

maar haar lichaam gebood
de gevoelde doodsnood.

Beklemming van de nek,
op haar borst een drukplek,

zweten in de handen,
het gevoel te landen

in de diepe afgrond,
haar droge open mond,

haar versnelde hartslag,
deden allen verslag

van de ervaren shock
die zich bij haar voltrok.

De sterke wil van haar
bleek onverzadigbaar

waardoor ze intens leed.
Wat zijn intrede deed

al na haar geboorte.
Haar in lijden loosde

en slingerde van pijn
naar een landerig zijn,

aldus Schopenhauer
als levensbeschouwer.

Had Emmy maar een pil
tegen haar taaie wil.

Emmy stond geestverstijfd.
Haar brein was ingelijfd

en wilde als het ging.
Dat gaf geruststelling.

“Kinderen die willen,
krijgen voor hun billen”,

was haar immers geleerd.
Ze had geaccepteerd

dat niemand de beker
voorbij kan laten gaan.

Niet zoals zij het wil,
maar zoals God het wil.

Volgens het humane
gebed in Getsemane.

“Ik wil dat ik het laat
anders dan het al gaat!

Ik spuug hier de vrucht uit,
die het goed en kwaad duidt.

Ik bezit niet de kennis
die alleen van God is.

De waan van het weten.
De waas van vergeten

dat ik slechts een mens ben.
Dat ik de weg niet ken.

Mijn bron van frustratie”,
aldus maalde Emmy.

Was het desastreuze
een bewuste keuze

of pure emotie?
Toonde dit de frictie

tussen haar hoofd en hart.
Nam het ontzag voor smart

werkelijk het besluit?
Kwam haar emotie uit

dat wat voelt ‘ik besta’?
Of was er anatta,

geen vast iets wat denkt
of ons gevoelens schenkt.

Anders kozen we wel
om nooit meer in de hel

te zijn als bestemming,
maar in een topstemming

ons leven te lijden.
De pijn te bestrijden.

Kon Emmy dit geschenk,
dat haar in een oogwenk

als een snel projectiel
volledig overviel,

toch bewust waarnemen
en een besluit nemen

op grond van harmonie?
Of was haar reactie

ook zonder vrije wil?
Gewaar zijn van een gril

zonder enige regie?
Dat gold niet voor Emmy.

Zij was geheel in shock
toen haar rede vertrok.

Van het bewogen ruige
was ze geen getuige.

“Je had niemand verwacht!
Niet aan een held gedacht.” -

Emmy zweeg - “Ik trotseer
altijd dit hondenweer

op jacht naar gevluchten.
Het zijn maar geruchten,

die we zelf verspreiden,
om jou te verleiden

te trappen in de schijn,
dat we bang en slap zijn.

Ik wachtte al uren
om jou te begluren.

Ik voorkwam ontdekking
door deze misvatting.”

De jager grapte blij:
“de waarheid ligt nabij,

daarom kijkt menigeen
er meestal overheen.”

De jager kon groeten
Emmy’s kleine voeten

met een plagende blik.
Hij had duidelijk schik.

Emmy’s spierwand spande.
Haar kringspier ontspande

en plaste langs benen
op haar kleine tenen.

De waarheid was haar shock.
Haar pis de aanwijsstok.

De jager werd hitsig.
Hij beval haar bitsig

haar bekken te draaien.
Hij wilde angst zaaien

zoals zijn vader deed.
Waaronder hij ooit leed

en bij haar herhaalde.
Hij het nu bepaalde.

Lust en dominantie
onderwierpen Emmy

alsof haar poseren
zou continueren

voor onbeperkte tijd.
Geen vergankelijkheid,

maar eeuwige paniek.
Blijvende erotiek.

Wat-als-scenario’s
ontsproten eindeloos

en waren levensecht.
Alsmaar werd ze geknecht.

Ze kende geen Sati
en dus verloor Emmy

contact met het heden,
en kwam het verleden

of de toekomst steeds op.
En met zo’n volle kop

gevuld met verhalen,
werd kalm ademhalen

een onmogelijkheid.
Ze raakte zichzelf kwijt

als reactie op vrees,
die zo opnieuw herrees.

Ze wilde van angst af,
zodat ze niet toegaf.

Ze zocht apatheia,
maar kreeg paranoia.

Als baby was angst ok.
Toen zat ze er niet mee.

Nu zag het als zwakte,
wat haar moed juist knakte

en slechts schaamte voortbracht
voor haar gebrek aan kracht.

“Wees flink. Stel je niet aan.
Niet meer bang zijn voortaan!”

Perfect zo aangeleerd
en geconditioneerd

door de samenleving.
De conditionering

stelde angst als vijand.
Dus groeide er geen band

van waaruit een begrip
en van daaruit wat grip

op waar de ‘foute’ vrees
eigenlijk naar verwees.

“Hier ben ik nog niet vrij.”
“Hiermee ben ik niet blij.”

Naar de onzekerheid
die haar brein zag als strijd.

Kon Emmy dit stoppen?
Zichzelf niet meer foppen

met valse oordelen?
Met haar angsten spelen?

Woelen zonder verhaal?
Het voelen zonder taal?

Zonder enige haast
wachten tot het uitblaast?

Ze kende niet de gril
van haar beperkte wil,

die haar frustratie bracht
als niet wat ze verwacht

in het echt was gebeurd
en het af had gekeurd.

En het dieper weten,
dat ze was vergeten,

dat wat haar wil wilde
vaak het feit niet stilde.

Dus werd het ego bang.
Zo nu en dan al lang

voordat het verwachtte
de voorliefde slachtte.

Emmy kon niet willen
wat mijn modegrillen

haar lieten ondergaan.
Zou het haar goed afgaan?

Haar Noche Oscura
start op de pagina,

nog even opgeschort,
tot de rijm duivels wordt.
